\lampiransection{Pedoman Wawancara Penelitian Utama}
\label{app:pedoman-wawancara-utama}

Lampiran ini digunakan sebagai pedoman wawancara penelitian utama. Pertanyaan disusun sebagai pohon pertanyaan agar peneliti dapat memulai dari pertanyaan pokok, lalu mencabangkan pertanyaan sesuai jawaban informan dan temuan awal observasi GrabFood.

\begin{flushleft}
\scriptsize
\begin{verbatim}
Akar: Strategi pemasaran Nasi Gerilya pada GrabFood
|-- Posisi kanal
|   |-- Peran GrabFood dalam penjualan
|   |-- Target pelanggan dan alasan memilih GrabFood
|   |-- Keterkaitan rating, ulasan, dan keputusan pembelian
|-- Produk dan nilai
|   |-- Menu unggulan, menu terlaris, porsi, dan rasa
|   |-- Menu tidak tersedia dan pengelolaan stok
|   |-- Perbandingan dengan menu kompetitor
|-- Harga dan promo
|   |-- Rentang harga, ongkir, voucher, dan margin
|   |-- Pengaruh promo terhadap pesanan
|   |-- Perbandingan harga kompetitor
|       |-- Paripurna, Dendeng Batokok, Istana Krakatau
|       |-- Pondok Gurih dan Garuda
|-- Proses layanan
|   |-- Alur pesanan, jam ramai, dan koordinasi dapur
|   |-- Akurasi pesanan, add-on, kemasan, dan komplain
|   |-- Perbaikan SOP layanan GrabFood
|-- Pelanggan dan persaingan
|   |-- Alasan membeli, membeli ulang, atau berhenti membeli
|   |-- Pengalaman yang tidak selalu muncul di review
|   |-- Pilihan kompetitor dan pembeda Nasi Gerilya
|-- Strategi
    |-- Kekuatan dan kelemahan internal
    |-- Peluang dan ancaman eksternal
    |-- Prioritas strategi peningkatan penjualan
\end{verbatim}
\end{flushleft}

\begin{table}[!htbp]
    \centering
    \caption{Pedoman Wawancara Pemilik atau Pengelola Berbasis Cabang Pertanyaan}
    \label{tab:pedoman-wawancara-pemilik-utama}
    \scriptsize
    \setlength{\tabcolsep}{2pt}
    \renewcommand{\arraystretch}{0.86}
    \begin{tabularx}{\textwidth}{|Z{2.2cm}|Y|Y|Y|}
        \toprule
        \textbf{Cabang} & \textbf{Pertanyaan Pokok} & \textbf{Cabang Lanjutan / Probing} & \textbf{Kaitan Observasi} \\ \hline
        \midrule
        Profil dan posisi kanal & Bagaimana perkembangan Nasi Gerilya dan mengapa GrabFood menjadi kanal penting? & Sejak kapan aktif di GrabFood? Seberapa besar kontribusinya terhadap penjualan? Kanal apa yang paling dibandingkan? & Menguji konteks TO-GF-01 dan klaim GrabFood sebagai kanal utama. \\ \hline
        STP & Siapa pelanggan utama yang ditargetkan melalui GrabFood dan citra apa yang ingin dibangun? & Apakah targetnya pekerja, keluarga, pelanggan sekitar, atau pelanggan promo? Apa citra utama: cepat, porsi besar, rasa Padang, atau harga? & Membantu menafsirkan rating 4,7, 4.584 penilaian, dan tema ulasan positif. \\ \hline
        Produk & Menu apa yang paling diunggulkan, paling sering dipesan, dan paling sering menimbulkan keluhan? & Mengapa menu terlaris tersebut kuat? Bagaimana standar porsi/rasa dijaga? Menu apa yang rawan kosong? & Menindaklanjuti TO-GF-04, TO-GF-06, TO-GF-07, dan TO-GF-08. \\ \hline
        Harga dan promo & Bagaimana harga, voucher, diskon, dan ongkir dipertimbangkan dalam strategi penjualan? & Apakah promo menaikkan volume pesanan? Bagaimana margin dijaga? Mengapa rentang harga dibuat lebar? & Menguji TO-GF-03 dan TO-GF-05; membandingkan TO-KP-06. \\ \hline
        Kompetitor & Siapa kompetitor utama dan aspek apa yang paling mengancam? & Apakah Paripurna, Dendeng Batokok, Istana Krakatau, Pondok Gurih, dan Garuda memang kompetitor relevan? Apakah ancamannya harga, rating, jumlah penilaian, jarak, reputasi legendaris, rasa, atau promo? & Menguji TO-KP-01--TO-KP-07. \\ \hline
        Operasional & Kendala apa yang paling sering muncul saat pesanan GrabFood ramai? & Bagaimana pembagian tugas admin dan dapur? Apa penyebab salah item, item kurang lengkap, atau add-on tidak diberikan? & Menindaklanjuti tema ulasan negatif/netral pada TO-GF-08. \\ \hline
        Pelanggan & Bagaimana kondisi pelanggan baru, pelanggan ulang, dan pelanggan yang berhenti membeli? & Apakah ada data repeat order? Apa alasan pelanggan kembali? Bagaimana komplain dicatat? & Mengurangi bias dari ulasan platform dan mendukung triangulasi pelanggan. \\ \hline
        Strategi & Strategi apa yang sudah dilakukan dan apa prioritas perbaikan berikutnya? & Jika harus memilih tiga prioritas, apakah produk, harga, promo, proses, tampilan toko, atau kompetitor? & Menjadi dasar IFAS, EFAS, dan SWOT pada Lampiran~\ref{app:template-ifas-efas-swot}. \\ \hline
        \bottomrule
    \end{tabularx}
    \tablesource{Diolah peneliti sebagai pedoman wawancara utama (2026).}
\end{table}

\begin{table}[!htbp]
    \centering
    \caption{Pedoman Wawancara Admin atau Karyawan Berbasis Cabang Pertanyaan}
    \label{tab:pedoman-wawancara-admin-karyawan}
    \scriptsize
    \setlength{\tabcolsep}{3pt}
    \renewcommand{\arraystretch}{0.92}
    \begin{tabularx}{\textwidth}{|Z{2.2cm}|Y|Y|Y|}
        \toprule
        \textbf{Cabang} & \textbf{Pertanyaan Pokok} & \textbf{Cabang Lanjutan / Probing} & \textbf{Kaitan Observasi} \\ \hline
        \midrule
        Alur pesanan & Bagaimana alur pesanan GrabFood diterima, diproses, dicek, dan diserahkan kepada driver? & Siapa yang menerima order? Siapa yang mengecek item? Apakah ada daftar cek sebelum diserahkan? & Menguji proses layanan yang berhubungan dengan TO-GF-08. \\ \hline
        Jam ramai & Pada jam atau hari apa pesanan paling ramai dan masalah apa yang sering muncul? & Apakah kesalahan lebih sering terjadi saat ramai? Bagaimana antrean driver ditangani? & Menghubungkan keluhan layanan dengan kondisi operasional aktual. \\ \hline
        Menu dan stok & Bagaimana pengelolaan menu habis atau stok yang belum diperbarui di aplikasi? & Siapa yang menutup menu di aplikasi? Berapa cepat status menu diperbarui? & Menindaklanjuti 6 menu tidak tersedia pada TO-GF-06. \\ \hline
        Akurasi pesanan & Apa penyebab pesanan tidak sesuai permintaan pelanggan? & Bagaimana add-on, kuah, sambal, lauk, dan catatan khusus dicek? & Menguji keluhan salah item, kelengkapan, dan add-on pada TO-GF-08. \\ \hline
        Kualitas produk & Bagaimana cara menjaga rasa, porsi, dan kondisi makanan saat order ramai? & Apakah ada standar porsi? Bagaimana makanan dipisah/dikemas? & Menguji tema rasa, porsi, kesegaran, dan kemasan pada TO-GF-07. \\ \hline
        Komplain & Keluhan apa yang paling sering diterima dan bagaimana tindak lanjutnya? & Apakah komplain dicatat? Apakah ada perbaikan setelah komplain berulang? & Membandingkan ulasan platform dengan pengalaman operasional. \\ \hline
        Perbaikan & Perbaikan apa yang paling dibutuhkan agar layanan GrabFood lebih lancar? & Apakah perlu SOP cek pesanan, foto menu baru, stok digital, atau koordinasi admin-dapur? & Menjadi bahan strategi proses dan \textit{physical evidence}. \\ \hline
        \bottomrule
    \end{tabularx}
    \tablesource{Diolah peneliti sebagai pedoman wawancara utama (2026).}
\end{table}

\begin{table}[!htbp]
    \centering
    \caption{Pedoman Wawancara Pelanggan Berbasis Cabang Pertanyaan}
    \label{tab:pedoman-wawancara-pelanggan}
    \scriptsize
    \setlength{\tabcolsep}{3pt}
    \renewcommand{\arraystretch}{0.92}
    \begin{tabularx}{\textwidth}{|Z{2.2cm}|Y|Y|Y|}
        \toprule
        \textbf{Cabang} & \textbf{Pertanyaan Pokok} & \textbf{Cabang Lanjutan / Probing} & \textbf{Kaitan Observasi} \\ \hline
        \midrule
        Pengalaman beli & Kapan terakhir membeli Nasi Gerilya melalui GrabFood dan menu apa yang dibeli? & Apakah membeli saat promo? Apakah pesanan untuk sendiri atau bersama? & Mengonfirmasi pengalaman aktual, bukan hanya persepsi umum. \\ \hline
        Alasan memilih & Apa alasan memilih Nasi Gerilya dibanding rumah makan Padang lain? & Apakah karena rasa, porsi, harga, promo, rating, jarak, atau kebiasaan? & Menguji makna rating dan ulasan positif pada TO-GF-01 dan TO-GF-07. \\ \hline
        Produk & Bagaimana penilaian terhadap rasa, porsi, kesegaran, variasi menu, dan foto menu? & Menu apa yang paling kuat? Menu apa yang mengecewakan? Apakah foto sesuai makanan? & Menghubungkan pengalaman pelanggan dengan TO-GF-04, TO-GF-07, dan TO-GF-08. \\ \hline
        Harga dan promo & Bagaimana penilaian terhadap harga, ongkir, voucher, dan promo Nasi Gerilya? & Pada harga berapa pelanggan merasa masih wajar? Apakah kompetitor terasa lebih murah? & Menguji TO-GF-03, TO-GF-05, dan TO-KP-06. \\ \hline
        Layanan & Bagaimana pengalaman terhadap kecepatan, akurasi pesanan, kemasan, dan kondisi makanan saat diterima? & Pernah ada salah item, item kurang, makanan kurang segar, atau catatan tidak dipenuhi? & Memeriksa tema keluhan TO-GF-08 dari sisi pelanggan. \\ \hline
        Review dan bias & Apakah pelanggan pernah mengalami masalah tetapi tidak menulis ulasan di GrabFood? & Jika tidak mengulas, mengapa? Apakah lebih sering mengulas ketika puas atau kecewa? & Menanggulangi \textit{survivorship bias} dalam pembacaan review. \\ \hline
        Kompetitor & Rumah makan Padang apa yang biasa dibandingkan dengan Nasi Gerilya dan mengapa? & Apakah pernah memilih Paripurna, Dendeng Batokok, Istana Krakatau, Pondok Gurih, Garuda, atau kompetitor lain karena rating, jumlah penilaian, harga, jarak, reputasi, atau menu? & Menguji relevansi TO-KP-01--TO-KP-07. \\ \hline
        Pembelian ulang & Apa yang membuat pelanggan mau atau tidak mau membeli ulang? & Apa satu perbaikan yang paling memengaruhi keputusan membeli lagi? & Menjadi dasar strategi retensi dan peningkatan penjualan. \\ \hline
        \bottomrule
    \end{tabularx}
    \tablesource{Diolah peneliti sebagai pedoman wawancara utama (2026).}
\end{table}

\clearpage
\subsection*{Format Pertanyaan dan Jawaban Wawancara}

Bagian berikut digunakan sebagai lembar kerja wawancara utama. Kolom jawaban diisi dengan kutipan penting atau ringkasan jawaban informan. Apabila informan memberi jawaban yang membuka tema baru, peneliti dapat menambahkan pertanyaan lanjutan pada kolom catatan lapangan.

\newcommand{\InterviewAnswerSpace}{\vspace{0.9cm}}
\newcommand{\InterviewLongtableCaption}[2]{%
    \refstepcounter{table}%
    \label{#1}%
    \addcontentsline{lot}{table}{\protect\numberline{\thetable}#2}%
    \begin{center}
        \footnotesize\bfseries
        Tabel \thetable\\
        #2
    \end{center}
}

\begingroup
\scriptsize
\setlength{\tabcolsep}{2pt}
\renewcommand{\arraystretch}{1.05}
\InterviewLongtableCaption{tab:tanya-jawab-wawancara-pemilik}{Tabel Pertanyaan dan Jawaban Wawancara Pemilik atau Pengelola}
\begin{longtable}{|>{\centering\arraybackslash}p{1cm}|>{\RaggedRight\arraybackslash}p{4.55cm}|>{\RaggedRight\arraybackslash}p{3.8cm}|>{\RaggedRight\arraybackslash}p{3.05cm}|}
    \hline
    \textbf{Kode} & \textbf{Pertanyaan Utama} & \textbf{Pertanyaan Pendalaman} & \textbf{Jawaban Informan} \\ \hline
    \endfirsthead
    \hline
    \textbf{Kode} & \textbf{Pertanyaan Utama} & \textbf{Pertanyaan Pendalaman} & \textbf{Jawaban Informan} \\ \hline
    \endhead
    W-PM-01 & Bagaimana sejarah singkat berdirinya Nasi Gerilya dan bagaimana perkembangan usaha sampai saat ini? & Sejak kapan beroperasi, siapa pengelola utama, perubahan menu/layanan apa yang paling penting, dan apa target usaha saat ini? & \InterviewAnswerSpace \\ \hline
    W-PM-02 & Mengapa GrabFood dipilih sebagai kanal utama dalam penjualan online? & Bandingkan dengan GoFood/ShopeeFood, kontribusi GrabFood terhadap penjualan, dan alasan kanal ini cocok dengan target pelanggan. & \InterviewAnswerSpace \\ \hline
    W-PM-03 & Siapa segmen pelanggan utama Nasi Gerilya di GrabFood? & Apakah pekerja, mahasiswa, keluarga, pelanggan sekitar, pelanggan promo, atau pelanggan rutin; bagaimana pola jam/hari pembelian mereka? & \InterviewAnswerSpace \\ \hline
    W-PM-04 & Citra atau posisi apa yang ingin dibangun Nasi Gerilya pada platform GrabFood? & Apakah ingin dikenal dari rasa Padang, porsi, harga, menu premium, kecepatan layanan, promo, atau reputasi/rating? & \InterviewAnswerSpace \\ \hline
    W-PM-05 & Menu apa yang paling diandalkan untuk menarik pesanan di GrabFood? & Sebutkan menu terlaris, alasan menu itu kuat, margin, ketersediaan stok, dan apakah sesuai dengan label terlaris pada observasi. & \InterviewAnswerSpace \\ \hline
    W-PM-06 & Bagaimana strategi harga Nasi Gerilya di GrabFood ditentukan? & Jelaskan pertimbangan biaya bahan, komisi platform, ongkir, promo, harga kompetitor, dan rentang harga Rp1.100--Rp275.000 yang terlihat pada observasi. & \InterviewAnswerSpace \\ \hline
    W-PM-07 & Bagaimana promo/voucher GrabFood digunakan dalam strategi pemasaran? & Apakah promo meningkatkan pesanan, bagaimana pengaruhnya pada margin, promo apa yang paling efektif, dan kapan promo biasanya digunakan? & \InterviewAnswerSpace \\ \hline
    W-PM-08 & Siapa kompetitor utama Nasi Gerilya di GrabFood menurut pemilik/pengelola? & Tanggapi Paripurna, Dendeng Batokok, Istana Krakatau, Pondok Gurih, Garuda, dan sebutkan kompetitor lain bila ada. & \InterviewAnswerSpace \\ \hline
    W-PM-09 & Aspek apa dari kompetitor yang paling mengancam Nasi Gerilya? & Bandingkan rating, jumlah penilaian, jarak, reputasi legendaris, harga bawah, harga premium, promo, foto menu, dan tema ulasan. & \InterviewAnswerSpace \\ \hline
    W-PM-10 & Masalah operasional apa yang paling sering memengaruhi pesanan GrabFood? & Jelaskan salah item, item kurang lengkap, add-on/kuah/sambal, menu kosong, kualitas makanan, jam ramai, dan koordinasi admin-dapur. & \InterviewAnswerSpace \\ \hline
    W-PM-11 & Bagaimana usaha menanggapi ulasan negatif atau komplain pelanggan? & Apakah komplain dicatat, siapa yang menindaklanjuti, contoh perbaikan yang pernah dilakukan, dan bagaimana pelanggan yang tidak menulis review dipahami. & \InterviewAnswerSpace \\ \hline
    W-PM-12 & Data kinerja apa yang dapat digunakan untuk menilai keberhasilan strategi pemasaran GrabFood? & Jumlah order, omzet, AOV, kunjungan toko, konversi, menu laris, pembelian ulang, komplain, periode promo, dan data sebelum-sesudah perbaikan. & \InterviewAnswerSpace \\ \hline
    W-PM-13 & Apa prioritas strategi peningkatan penjualan Nasi Gerilya dalam waktu dekat? & Urutkan prioritas produk, harga, promo, tampilan toko, SOP pesanan, stok digital, foto menu, rating/ulasan, dan pembeda dari kompetitor. & \InterviewAnswerSpace \\ \hline
\end{longtable}
\addtocounter{table}{-1}
\tablesource[\textwidth]{Diolah peneliti sebagai format wawancara pemilik atau pengelola penelitian utama (2026).}
\endgroup

\begingroup
\scriptsize
\setlength{\tabcolsep}{2pt}
\renewcommand{\arraystretch}{1.05}
\InterviewLongtableCaption{tab:tanya-jawab-wawancara-admin}{Tabel Pertanyaan dan Jawaban Wawancara Admin atau Karyawan}
\begin{longtable}{|>{\centering\arraybackslash}p{1cm}|>{\RaggedRight\arraybackslash}p{4.55cm}|>{\RaggedRight\arraybackslash}p{3.8cm}|>{\RaggedRight\arraybackslash}p{3.05cm}|}
    \hline
    \textbf{Kode} & \textbf{Pertanyaan Utama} & \textbf{Pertanyaan Pendalaman} & \textbf{Jawaban Informan} \\ \hline
    \endfirsthead
    \hline
    \textbf{Kode} & \textbf{Pertanyaan Utama} & \textbf{Pertanyaan Pendalaman} & \textbf{Jawaban Informan} \\ \hline
    \endhead
    W-AD-01 & Bagaimana alur pesanan GrabFood diterima sampai diserahkan kepada driver? & Siapa menerima order, siapa menyiapkan makanan, siapa mengecek kelengkapan, dan apakah ada pembagian tugas tertulis? & \InterviewAnswerSpace \\ \hline
    W-AD-02 & Pada jam atau hari apa pesanan GrabFood biasanya paling ramai? & Apa dampak jam ramai terhadap kecepatan, akurasi, antrean driver, stok lauk, dan koordinasi dapur? & \InterviewAnswerSpace \\ \hline
    W-AD-03 & Bagaimana cara memastikan pesanan sesuai dengan catatan pelanggan? & Bagaimana add-on, kuah, sambal, lauk, minuman, dan catatan khusus dicek sebelum pesanan keluar? & \InterviewAnswerSpace \\ \hline
    W-AD-04 & Kesalahan pesanan apa yang paling sering terjadi? & Salah item, item kurang, add-on tidak masuk, pesanan tertukar, menu habis, keterlambatan, atau makanan kurang sesuai harapan. & \InterviewAnswerSpace \\ \hline
    W-AD-05 & Bagaimana pengelolaan menu yang habis atau tidak tersedia di aplikasi? & Siapa yang menutup menu, kapan status diperbarui, apa penyebab 6 item tidak tersedia, dan bagaimana dampaknya pada pelanggan? & \InterviewAnswerSpace \\ \hline
    W-AD-06 & Bagaimana standar porsi, rasa, dan pengemasan dijaga saat order ramai? & Apakah ada ukuran porsi, standar lauk, standar sambal/kuah, pemisahan makanan, dan pengecekan kondisi sebelum dikirim? & \InterviewAnswerSpace \\ \hline
    W-AD-07 & Keluhan pelanggan apa yang paling sering diterima dari pesanan GrabFood? & Bandingkan keluhan langsung dengan ulasan platform: rasa, porsi, kelengkapan, kualitas, harga, atau pelayanan. & \InterviewAnswerSpace \\ \hline
    W-AD-08 & Bagaimana tindak lanjut jika ada komplain atau rating buruk? & Siapa yang merespons, apakah ada kompensasi, apakah masalah dicatat, dan apakah ada perubahan SOP setelah komplain berulang? & \InterviewAnswerSpace \\ \hline
    W-AD-09 & Bagaimana komunikasi antara admin, dapur, kasir, dan driver dilakukan? & Apakah menggunakan catatan, aplikasi, panggilan langsung, atau pengecekan lisan; titik rawan miskomunikasi di mana? & \InterviewAnswerSpace \\ \hline
    W-AD-10 & Apa perbaikan paling realistis agar pesanan GrabFood lebih akurat dan cepat? & SOP cek pesanan, daftar cek, pembaruan stok, foto menu, tambahan staf saat ramai, pelabelan pesanan, atau pelatihan. & \InterviewAnswerSpace \\ \hline
    W-AD-11 & Data operasional apa yang tersedia untuk mendukung analisis? & Catatan komplain, menu habis, pesanan batal, jam ramai, menu paling sering dipesan, dan kendala harian. & \InterviewAnswerSpace \\ \hline
\end{longtable}
\addtocounter{table}{-1}
\tablesource[\textwidth]{Diolah peneliti sebagai format wawancara admin atau karyawan penelitian utama (2026).}
\endgroup

\begingroup
\scriptsize
\setlength{\tabcolsep}{2pt}
\renewcommand{\arraystretch}{1.05}
\InterviewLongtableCaption{tab:tanya-jawab-wawancara-kompetitor}{Tabel Pertanyaan dan Jawaban Wawancara Kompetitor atau Usaha Sejenis}
\begin{longtable}{|>{\centering\arraybackslash}p{1cm}|>{\RaggedRight\arraybackslash}p{4.55cm}|>{\RaggedRight\arraybackslash}p{3.8cm}|>{\RaggedRight\arraybackslash}p{3.05cm}|}
    \hline
    \textbf{Kode} & \textbf{Pertanyaan Utama} & \textbf{Pertanyaan Pendalaman} & \textbf{Jawaban Informan} \\ \hline
    \endfirsthead
    \hline
    \textbf{Kode} & \textbf{Pertanyaan Utama} & \textbf{Pertanyaan Pendalaman} & \textbf{Jawaban Informan} \\ \hline
    \endhead
    W-KP-01 & Bagaimana profil singkat usaha dan sejak kapan menggunakan platform pesan antar seperti GrabFood? & Nama usaha, lokasi, jenis menu utama, kanal online yang digunakan, dan peran GrabFood dalam penjualan. & \InterviewAnswerSpace \\ \hline
    W-KP-02 & Siapa pelanggan utama usaha ini pada platform pesan antar? & Pelanggan sekitar, pekerja, keluarga, pelanggan promo, pelanggan rutin, atau pelanggan yang mencari menu tertentu. & \InterviewAnswerSpace \\ \hline
    W-KP-03 & Faktor apa yang paling penting agar pelanggan memilih usaha ini di GrabFood? & Rasa, harga, porsi, promo, rating, ulasan, jarak, foto menu, reputasi lama, atau kelengkapan menu. & \InterviewAnswerSpace \\ \hline
    W-KP-04 & Bagaimana strategi harga dan paket menu ditentukan? & Pertimbangan komisi platform, harga bahan, porsi, menu premium, harga bawah, promo, dan perbandingan dengan usaha sejenis. & \InterviewAnswerSpace \\ \hline
    W-KP-05 & Bagaimana promo atau voucher digunakan untuk menarik pelanggan? & Jenis promo yang sering dipakai, pengaruh terhadap order, risiko margin, dan waktu paling tepat menggunakan promo. & \InterviewAnswerSpace \\ \hline
    W-KP-06 & Bagaimana usaha menjaga rating dan ulasan pelanggan di platform? & Tindakan setelah komplain, cara menjaga kualitas, apakah rating dipantau, dan bagaimana respon terhadap ulasan negatif. & \InterviewAnswerSpace \\ \hline
    W-KP-07 & Bagaimana usaha mengelola ketersediaan menu dan stok pada aplikasi? & Siapa yang memperbarui stok, kapan menu ditutup, item apa yang paling rawan habis, dan dampaknya pada rating/order. & \InterviewAnswerSpace \\ \hline
    W-KP-08 & Kendala operasional apa yang paling sering terjadi dalam pesanan online? & Salah item, keterlambatan, antrean driver, kemasan, menu habis, catatan pelanggan, atau kualitas makanan setelah pengiriman. & \InterviewAnswerSpace \\ \hline
    W-KP-09 & Menurut informan, apa yang membedakan usaha ini dari rumah makan Padang lain di GrabFood? & Menu khas, reputasi, rasa, porsi, harga, kecepatan, layanan, jam buka, foto menu, atau lokasi. & \InterviewAnswerSpace \\ \hline
    W-KP-10 & Bagaimana informan melihat persaingan rumah makan Padang di GrabFood area Medan? & Seberapa ketat persaingan, faktor penentu kemenangan, dan apakah pelanggan mudah berpindah ke merchant lain. & \InterviewAnswerSpace \\ \hline
    W-KP-11 & Apa saran umum bagi rumah makan Padang agar penjualan di GrabFood meningkat? & Perbaikan produk, harga, promo, foto, SOP, stok, rating, ulasan, kemasan, atau menu unggulan. & \InterviewAnswerSpace \\ \hline
\end{longtable}
\addtocounter{table}{-1}
\tablesource[\textwidth]{Diolah peneliti sebagai format wawancara kompetitor atau usaha sejenis penelitian utama (2026).}
\endgroup

\begingroup
\scriptsize
\setlength{\tabcolsep}{2pt}
\renewcommand{\arraystretch}{1.05}
\InterviewLongtableCaption{tab:tanya-jawab-wawancara-konsumen}{Tabel Pertanyaan dan Jawaban Wawancara Konsumen}
\begin{longtable}{|>{\centering\arraybackslash}p{1cm}|>{\RaggedRight\arraybackslash}p{4.55cm}|>{\RaggedRight\arraybackslash}p{3.8cm}|>{\RaggedRight\arraybackslash}p{3.05cm}|}
    \hline
    \textbf{Kode} & \textbf{Pertanyaan Utama} & \textbf{Pertanyaan Pendalaman} & \textbf{Jawaban Informan} \\ \hline
    \endfirsthead
    \hline
    \textbf{Kode} & \textbf{Pertanyaan Utama} & \textbf{Pertanyaan Pendalaman} & \textbf{Jawaban Informan} \\ \hline
    \endhead
    W-PL-01 & Kapan terakhir membeli Nasi Gerilya melalui GrabFood dan menu apa yang dibeli? & Frekuensi pembelian, menu yang paling sering dibeli, membeli untuk sendiri/keluarga/kantor, dan waktu pembelian. & \InterviewAnswerSpace \\ \hline
    W-PL-02 & Apa alasan utama memilih Nasi Gerilya dibanding rumah makan Padang lain? & Rasa, porsi, harga, promo, rating, ulasan, jarak, kebiasaan, menu tertentu, atau rekomendasi orang lain. & \InterviewAnswerSpace \\ \hline
    W-PL-03 & Bagaimana penilaian terhadap rasa dan konsistensi makanan Nasi Gerilya? & Menu yang paling enak, menu yang kurang sesuai, perubahan rasa, kesegaran bahan, dan perbandingan dengan kompetitor. & \InterviewAnswerSpace \\ \hline
    W-PL-04 & Bagaimana penilaian terhadap porsi, lauk, sambal, kuah, dan kelengkapan pesanan? & Apakah sesuai harga, pernah ada item kurang, add-on tidak masuk, sambal/kuah kurang, atau porsi tidak konsisten. & \InterviewAnswerSpace \\ \hline
    W-PL-05 & Bagaimana penilaian terhadap harga, ongkir, voucher, dan promo Nasi Gerilya? & Apakah terasa murah/setara/mahal, promo memengaruhi pembelian, dan harga berapa yang dianggap wajar. & \InterviewAnswerSpace \\ \hline
    W-PL-06 & Bagaimana penilaian terhadap foto menu dan tampilan toko di GrabFood? & Apakah foto sesuai makanan, menu mudah dipilih, informasi harga jelas, dan tampilan toko membuat percaya untuk membeli. & \InterviewAnswerSpace \\ \hline
    W-PL-07 & Bagaimana pengalaman layanan saat pesanan diterima? & Kecepatan, kemasan, suhu makanan, kondisi makanan, ketepatan catatan, dan pengalaman dengan driver/aplikasi. & \InterviewAnswerSpace \\ \hline
    W-PL-08 & Apakah pernah mengalami masalah saat membeli Nasi Gerilya melalui GrabFood? & Salah item, item kurang, makanan kurang segar, rasa berubah, menu habis, harga/porsi tidak sesuai, atau keterlambatan. & \InterviewAnswerSpace \\ \hline
    W-PL-09 & Jika pernah mengalami masalah, apakah menulis ulasan di GrabFood? & Jika tidak menulis, alasannya apa; apakah lebih sering mengulas saat puas, kecewa, atau tidak pernah mengulas. & \InterviewAnswerSpace \\ \hline
    W-PL-10 & Rumah makan Padang apa yang biasa dibandingkan dengan Nasi Gerilya? & Paripurna, Dendeng Batokok, Istana Krakatau, Pondok Gurih, Garuda, atau merchant lain; apa alasan membandingkan. & \InterviewAnswerSpace \\ \hline
    W-PL-11 & Apa yang membuat konsumen memilih kompetitor dibanding Nasi Gerilya? & Harga lebih murah, rating lebih tinggi, jarak lebih dekat, promo lebih menarik, menu tertentu, reputasi, atau pengalaman sebelumnya. & \InterviewAnswerSpace \\ \hline
    W-PL-12 & Apa yang membuat konsumen mau membeli ulang Nasi Gerilya? & Rasa, porsi, promo, konsistensi, kelengkapan pesanan, harga, rating, menu favorit, atau pengalaman layanan. & \InterviewAnswerSpace \\ \hline
    W-PL-13 & Perbaikan apa yang paling diharapkan dari Nasi Gerilya di GrabFood? & Produk, harga, promo, foto menu, pilihan menu, stok, kemasan, kecepatan, akurasi pesanan, atau respon komplain. & \InterviewAnswerSpace \\ \hline
\end{longtable}
\addtocounter{table}{-1}
\tablesource[\textwidth]{Diolah peneliti sebagai format wawancara konsumen penelitian utama (2026).}
\endgroup
